\documentclass[11pt,a4paper]{article}
\usepackage[T1]{fontenc}
\usepackage{lmodern}
\usepackage[margin=27mm]{geometry}
\usepackage{amsmath,amssymb,amsthm,mathtools}
\usepackage{microtype}
\usepackage{booktabs,array}
\usepackage{enumitem}
\usepackage[hidelinks]{hyperref}
\setlength{\parskip}{5pt plus 1pt}
\setlength{\parindent}{0pt}
\setlength{\emergencystretch}{3em}
\setlist{itemsep=2pt,topsep=3pt}
\raggedbottom
\newtheorem{theorem}{Theorem}[section]
\newtheorem{lemma}[theorem]{Lemma}
\newtheorem{corollary}[theorem]{Corollary}
\newtheorem{proposition}[theorem]{Proposition}
\theoremstyle{definition}
\newtheorem{definition}[theorem]{Definition}
\newtheorem{example}[theorem]{Example}
\theoremstyle{remark}
\newtheorem{remark}[theorem]{Remark}
\newcommand{\Z}{\mathbb{Z}}
\newcommand{\N}{\mathbb{N}}
\newcommand{\Q}{\mathbb{Q}}
\newcommand{\F}{\mathbb{F}}
\newcommand{\ps}[1]{#1[[x]]}
\newcommand{\coeff}[2]{[x^{#1}]\,#2}
\newcommand{\hl}[1]{\textbf{#1}}
\newcommand{\para}[1]{\par\medskip\noindent\textbf{#1}\enspace}
\hypersetup{pdftitle={A power series whose square is x² + 2x³ modulo 9},pdfauthor={Dev D. Goyal},pdfsubject={OEIS A389474}}
\title{A power series whose square is $x^2+2x^3$ modulo $9$\\[6pt]\large A proof of two conjectures in OEIS A389474}
\author{Dev D. Goyal}
\date{October 2026}
\begin{document}
\maketitle
\begin{abstract}
The sequence A389474 of the On-Line Encyclopedia of Integer Sequences (OEIS) lists the coefficients of the power series
$A(x)=x+x^2+4x^3+14x^4+\cdots$ that satisfies $A\bigl(x^2+2x^3+2A(x)^2\bigr)=3A(x)^2$. Paul~D.~Hanna conjectured that
all coefficients of the square $A(x)^2$ beyond $x^3$ are divisible by $9$, and that every coefficient of
$A(x)^{3k}-A(x^3)^k$ is divisible by $3$. We prove both. The main result is that $A(x)^2\equiv x^2+2x^3\pmod 9$. We also prove that the
equation has exactly one normalized solution with integer coefficients; this is not immediate, because computing the
coefficients from the equation involves a division by $2$ at every step. The proof works with the inverse series of $A$, which satisfies a simpler equation. The second
conjecture is an instance of a classical identity. The results have been checked by computer and proved formally in
the Lean proof assistant. A short further argument, not part of the formal verification, shows that $A(x)^2$ is
congruent to a polynomial modulo every power of $3$.
\end{abstract}
\begin{center}\begin{minipage}{0.88\linewidth}\small
\textbf{How this paper was produced.} This paper reports the result of an AI research effort directed by
Dev~D.~Goyal. ChatGPT (OpenAI), using the GPT-6 model with Pro reasoning and working autonomously from prompts by
Dev~D.~Goyal, found the proof, wrote its Lean~4 formalization, and wrote a first draft of the paper.
Dev~D.~Goyal chose the problems and asked for proofs, formal verification and revisions, but did not supply proof
ideas. Claude Opus~5.5, used with Claude Code, then checked the work independently: it rebuilt the formal proofs from
source and confirmed that they rely only on Lean's three standard axioms, recomputed the sequence from its definition
compared it with the published terms and tested every claim numerically, and read every step of the proofs. It also wrote the present
version of this paper.
Dev~D.~Goyal has verified the results and is responsible for them. The paper has not been peer reviewed.
The section ``About this paper'' at the end gives the details.
\end{minipage}\end{center}

\section{Introduction}

\subsection{The sequence}\label{sec:sequence}
The \emph{On-Line Encyclopedia of Integer Sequences} (OEIS)~\cite{oeisdb} is a large public database of integer
sequences. Many
entries also record \emph{conjectures}: patterns that the contributor observed in the data but did not prove.

This paper is about entry A389474~\cite{oeis}, contributed by Paul~D.~Hanna in October 2025. It defines a power series
$A(x)=a(1)\,x+a(2)\,x^2+a(3)\,x^3+\cdots$ with integer coefficients and $a(1)=1$ by the equation
\begin{equation}\label{eq:oeis}
A\bigl(x^2+2x^3+2A(x)^2\bigr)=3A(x)^2 .
\end{equation}
The left side means: substitute the series $x^2+2x^3+2A(x)^2$ into $A$ (Section~\ref{sec:background} explains
substitution). On the right, $A(x)^2=A(x)\cdot A(x)$ is an ordinary square. Theorem~\ref{thm:exist} shows that there
is exactly one such series. Its first coefficients, starting at $n=1$, are 1, 1, 4, 14, 104, 488, 3291, 18141, 127527, 790427, \dots{}
The entry reads, in part:
\begin{quote}\small\ttfamily\raggedright G.f. A(x) satisfies A( x\^{}2 + 2*x\^{}3 + 2*A(x)\^{}2 ) = 3*A(x)\^{}2.\par\smallskip (2) A( 2*x\^{}2 + R(x)\^{}2 + 2*R(x)\^{}3 ) = 3*x\^{}2 where A(R(x)) = x.\par\smallskip (3) [x\^{}n] A(x)\^{}2 == 0 (mod 9) for n > 3 (conjecture).\par\smallskip (4) [x\^{}n] A(x)\^{}(3*k) - A(x\^{}3)\^{}k == 0 (mod 3) for n >= 1 and k >= 1 (conjecture).\par\smallskip\hfill\textrm{--- OEIS A389474, Paul D. Hanna, October 2025}\end{quote}

Here ``G.f.''\ stands for \emph{generating function}, the series whose coefficients are the terms of the sequence, and
$[x^n]\,F$ denotes the coefficient of $x^n$ in a series $F$. To avoid confusion with the equation numbers of this paper,
we refer to the formulas numbered (2), (3) and (4) in the entry as (E2), (E3) and (E4). In (E2), $R$ is the
compositional inverse of $A$: the series that undoes the substitution of $A$ (Section~\ref{sec:background}).

\subsection{The question}
Statement (E3) says that the coefficients of the square $A(x)^2$ are divisible by $9$ from $x^4$ on. Statement
(E4) says that for every $k\ge1$ all coefficients of $A(x)^{3k}-A(x^3)^k$ are divisible by $3$; here $A(x)^{3k}$ is a
power of $A$, and $A(x^3)$ is $A$ with $x$ replaced by $x^3$. The table shows the coefficients of $A(x)^2$ modulo $9$
(the first terms of $A(x)^2$ are $x^2+2x^3+9x^4+36x^5+252x^6+\cdots$). The only nonzero entries, at $n=2$ and $n=3$,
are printed in bold.
\begin{center}\small\setlength{\tabcolsep}{3.2pt}
\begin{tabular}{>{\raggedright\arraybackslash}p{7.4em}>{\centering\arraybackslash}p{1.15em}>{\centering\arraybackslash}p{1.15em}>{\centering\arraybackslash}p{1.15em}>{\centering\arraybackslash}p{1.15em}>{\centering\arraybackslash}p{1.15em}>{\centering\arraybackslash}p{1.15em}>{\centering\arraybackslash}p{1.15em}>{\centering\arraybackslash}p{1.15em}>{\centering\arraybackslash}p{1.15em}>{\centering\arraybackslash}p{1.15em}>{\centering\arraybackslash}p{1.15em}>{\centering\arraybackslash}p{1.15em}>{\centering\arraybackslash}p{1.15em}>{\centering\arraybackslash}p{1.15em}>{\centering\arraybackslash}p{1.15em}>{\centering\arraybackslash}p{1.15em}>{\centering\arraybackslash}p{1.15em}>{\centering\arraybackslash}p{1.15em}>{\centering\arraybackslash}p{1.15em}>{\centering\arraybackslash}p{1.15em}}
\toprule
$n$ & 0 & 1 & \hl{2} & \hl{3} & 4 & 5 & 6 & 7 & 8 & 9 & 10 & 11 & 12 & 13 & 14 & 15 & 16 & 17 & 18 & 19 \\ \midrule
$\coeff{n}{A^2}\bmod 9$ & 0 & 0 & \hl{1} & \hl{2} & 0 & 0 & 0 & 0 & 0 & 0 & 0 & 0 & 0 & 0 & 0 & 0 & 0 & 0 & 0 & 0 \\
\bottomrule
\end{tabular}
\\[4pt]
\begin{tabular}{>{\raggedright\arraybackslash}p{7.4em}>{\centering\arraybackslash}p{1.15em}>{\centering\arraybackslash}p{1.15em}>{\centering\arraybackslash}p{1.15em}>{\centering\arraybackslash}p{1.15em}>{\centering\arraybackslash}p{1.15em}>{\centering\arraybackslash}p{1.15em}>{\centering\arraybackslash}p{1.15em}>{\centering\arraybackslash}p{1.15em}>{\centering\arraybackslash}p{1.15em}>{\centering\arraybackslash}p{1.15em}>{\centering\arraybackslash}p{1.15em}>{\centering\arraybackslash}p{1.15em}>{\centering\arraybackslash}p{1.15em}>{\centering\arraybackslash}p{1.15em}>{\centering\arraybackslash}p{1.15em}>{\centering\arraybackslash}p{1.15em}>{\centering\arraybackslash}p{1.15em}>{\centering\arraybackslash}p{1.15em}>{\centering\arraybackslash}p{1.15em}>{\centering\arraybackslash}p{1.15em}}
\toprule
$n$ & 20 & 21 & 22 & 23 & 24 & 25 & 26 & 27 & 28 & 29 & 30 & 31 & 32 & 33 & 34 & 35 & 36 & 37 & 38 & 39 \\ \midrule
$\coeff{n}{A^2}\bmod 9$ & 0 & 0 & 0 & 0 & 0 & 0 & 0 & 0 & 0 & 0 & 0 & 0 & 0 & 0 & 0 & 0 & 0 & 0 & 0 & 0 \\
\bottomrule
\end{tabular}
\end{center}

\subsection{What this paper proves}
\textbf{Main result} (Theorem~\ref{thm:nine}). \emph{Modulo $9$, the square of $A$ is a polynomial:}
\[
A(x)^2\equiv x^2+2x^3\pmod 9 .
\]
So conjecture (E3) is true. Conjecture (E4) is true as well (Proposition~\ref{prop:frob}); in fact it holds for
\emph{every} power series with integer coefficients, by a classical identity. We also prove that \eqref{eq:oeis} has
exactly one normalized solution with integer coefficients (Theorem~\ref{thm:exist}). This needs an argument: when the coefficients
are computed from \eqref{eq:oeis} one at a time, each new coefficient $a(n)$ first appears in degree $n+1$ with the
factor $-2$, so each step involves a division by $2$.

The key idea is to pass to the inverse series. Applying $R$ to both sides of \eqref{eq:oeis} gives
$x^2+2x^3+2A(x)^2=R\bigl(3A(x)^2\bigr)$. Modulo $9$, the square of $3A(x)^2$ is $9A(x)^4\equiv0$, so all terms of $R$
beyond the first contribute nothing, and $R\bigl(3A(x)^2\bigr)\equiv3A(x)^2$. The equation then says
$x^2+2x^3+2A(x)^2\equiv3A(x)^2$, that is, $A(x)^2\equiv x^2+2x^3$.

Section~\ref{sec:discussion} pushes the same idea further: $A(x)^2$ is congruent to a polynomial modulo every power
of $3$, and $A$ itself is determined modulo $9$ by an explicit formula.

\subsection{What was known}
The entry records the equation (E2) satisfied by the inverse series $R$, and statements (E3) and (E4) as
conjectures. Hanna also contributed the companion entry A389475~\cite{oeisother}, whose generating function satisfies a
closely related equation; it is treated in a companion paper with the same method. We are not aware of an earlier
proof of conjecture (E3). Statement (E4) follows directly from the Frobenius identity (Lemma~\ref{lem:frobenius}), as
shown in Proposition~\ref{prop:frob}; we include the two-line proof for completeness and claim nothing new about it.

\subsection{Motivation}
Equations in which an unknown series is substituted into an expression built from itself are common in
combinatorics. The entry itself compares \eqref{eq:oeis} with the identity $C\bigl(x^2+2x\,C(x)^2\bigr)=C(x)^2$ for the
generating series $C$ of the Catalan numbers. The coefficients of $A$ grow quickly and look irregular, yet modulo $9$
the square of $A$ is a polynomial: everything beyond $x^3$ vanishes. The proof uses two techniques with wider application.
First, an equation that seems to require division by $2$ can still have an integer solution, if the numbers being
divided are always even; here this follows from an identity modulo $2$. Second, replacing a series by its inverse
can reduce an equation to one that can be solved explicitly modulo a power of a prime.

\subsection{How to read this paper}
Section~\ref{sec:background} explains the notions we need: power series, substitution, compositional inverses, and
remainders. Section~\ref{sec:results} states the results, and Section~\ref{sec:proofs} proves them.
Section~\ref{sec:checks} describes how the results were checked, by computer and by a formal proof in Lean.
Section~\ref{sec:discussion} discusses the scope of the results, some stronger statements, and open questions.
\section{Background}\label{sec:background}
This section collects the definitions and facts used later. Readers who know formal power series and congruences can skip it and return when needed.
\subsection{Formal power series}
A \emph{formal power series} is an expression
\[
U(x)=u_0+u_1x+u_2x^2+u_3x^3+\cdots ,
\]
determined by its sequence of coefficients $u_0,u_1,u_2,\dots$; see~\cite{niven,stanley}. No number is ever
substituted for $x$, so questions of convergence do not arise. We write $\coeff{n}{U}=u_n$ for the coefficient of $x^n$ in $U$; two series are equal
when all their coefficients are equal.

Series are added coefficient by coefficient and multiplied like polynomials:
\[
\coeff{n}{(UV)}=u_0v_n+u_1v_{n-1}+\cdots+u_nv_0 .
\]
Each coefficient of a product is a finite sum. For example, $(1-x)(1+x+x^2+x^3+\cdots)=1$, so $1+x+x^2+\cdots$ is
the \emph{multiplicative inverse} of $1-x$, written $\frac{1}{1-x}$. More generally, every series $D$ with constant term
$1$ has a unique multiplicative inverse $\frac1D$, whose coefficients are determined one at a time by $D\cdot\frac1D=1$
using only addition, subtraction and multiplication. Whenever a quotient of power series appears in this paper, its
denominator has constant term $1$, so the quotient is again a power series.

A \emph{commutative ring} is a set with an addition, a subtraction and a multiplication that satisfy the usual rules
of arithmetic; examples are $\Z$, $\Q$ and the integers modulo $m$. For a commutative ring $R$ we write $\ps{R}$ for
the set of power series with coefficients in $R$; for instance, writing $\F_2$ for the integers modulo $2$, $\ps{\F_2}$
consists of the series with coefficients modulo $2$. We write
$U(0)$ for the constant term $u_0$ of $U$, and $U=O(x^d)$ when the coefficients of $1,x,\dots,x^{d-1}$ in $U$ are all
zero. Thus $U-V=O(x^d)$ means that $U$ and $V$ have the same coefficient of $x^n$ for every $n<d$. A series with
constant term $0$ and coefficient of $x$ equal to $1$, that is, a series of the form $x+u_2x^2+u_3x^3+\cdots$, is
called \emph{normalized}.

\subsection{Substituting one series into another}
If $V(x)$ has constant term $0$, we can substitute $V$ for $x$ in $U$:
\[
U(V(x))=u_0+u_1V(x)+u_2V(x)^2+u_3V(x)^3+\cdots .
\]
This is called \emph{composition}, also written $U\circ V$. It makes sense because $V(x)^k$ has no terms below $x^k$,
so only the finitely many terms with $k\le n$ affect the coefficient of $x^n$. For example, if $U=x+x^2$ and $V=x^2$, then
$U(V(x))=x^2+x^4$. Composition is associative: $U\circ(V\circ W)=(U\circ V)\circ W$.

Substituting a series into itself repeatedly gives its \emph{iterates}: $U^{\circ2}(x)=U(U(x))$,
$U^{\circ3}(x)=U(U(U(x)))$, and so on, with $U^{\circ0}(x)=x$ and $U^{\circ1}=U$. They must not be confused with
\emph{powers} $U^2=U\cdot U$, $U^3=U\cdot U\cdot U$. (The OEIS sometimes writes $A^k(x)$ for the $k$-th iterate. In
this paper the small circle always means an iterate, and a plain exponent always means a power.)

\subsection{Undoing a substitution}
A normalized series can always be ``undone'' by another normalized series.

\begin{lemma}[Compositional inverse]\label{lem:inverse}
For every normalized series $U$ there is exactly one normalized series $V$ with $U(V(x))=x$, and this $V$ also
satisfies $V(U(x))=x$. If $U$ has integer coefficients, so does $V$.
\end{lemma}

We call $V$ the \emph{compositional inverse} of $U$. The OEIS calls it the \emph{series reversion} of $U$. For example,
the compositional inverse of $x-x^2$ is $C(x)=x+x^2+2x^3+5x^4+14x^5+\cdots$, whose coefficients are the Catalan numbers.
Indeed $C-C^2=x$, which says exactly that substituting $C$ into $x-x^2$ gives $x$.

\begin{proof}
Write $V=x+v_2x^2+v_3x^3+\cdots$ with unknown coefficients. In
$U(V)=V+u_2V^2+u_3V^3+\cdots$, the coefficient of $x^n$ (for $n\ge2$) equals $v_n$ plus terms that involve only
$v_2,\dots,v_{n-1}$. Setting it equal to $0$ therefore determines $v_n$ from the earlier coefficients, using only
additions, subtractions and multiplications. So there is exactly one $V$ with $U\circ V=x$, and it has integer
coefficients if $U$ does. Applying the same construction to $V$ gives a normalized $W$ with $V\circ W=x$. Then
$U=U\circ(V\circ W)=(U\circ V)\circ W=W$, so $V\circ U=V\circ W=x$ as well.
\end{proof}

\subsection{Remainders and congruences}
For an integer $m\ge2$ we write $a\equiv b\pmod m$ (``$a$ is congruent to $b$ modulo $m$'') when $a-b$ is divisible by
$m$, that is, when $a$ and $b$ leave the same remainder after division by $m$. The remainder of $a$ on division by
$m$ is called its \emph{residue} modulo $m$; it is one of $0,1,\dots,m-1$.

Addition and multiplication are compatible with congruences, so the residues $0,1,\dots,m-1$, added and multiplied
modulo $m$, form a ring, written $\Z/m\Z$. When $m=p$ is a prime we also write $\F_p$. In $\F_2=\{0,1\}$ we have $1+1=0$,
so addition and subtraction coincide.

For power series with integer coefficients, $U\equiv V\pmod m$ means that \emph{every} coefficient of $U-V$ is
divisible by $m$. Equivalently, the two series become equal when every coefficient is replaced by its residue; we call
this \emph{reducing} the series modulo $m$. The key fact is the following.

\emph{Reduction modulo $m$ is compatible with all the operations above.} To reduce a sum, a product, a substitution, or
an inverse, we may reduce the ingredients first and then compute in $\Z/m\Z$. This holds because every
coefficient of the result is obtained from finitely many coefficients of the ingredients by adding, subtracting and
multiplying, and these operations are compatible with remainders.

\para{Rings.} Some statements below hold over every commutative ring, for example $\Z$, $\Q$ and $\Z/m\Z$. No
division is used in them; this matters because in $\Z/m\Z$ division is not always possible (for instance, $2$ has no inverse in
$\Z/4\Z$).

\subsection{The Frobenius identity}
\begin{lemma}[Frobenius identity]\label{lem:frobenius}
Let $p$ be a prime. For every power series $U$ with coefficients in $\F_p$, we have $U(x)^p=U(x^p)$.
\end{lemma}

\begin{proof}
In $\F_p$, the binomial theorem simplifies to $(a+b)^p=a^p+b^p$, because every binomial coefficient
$\binom pi$ with $0<i<p$ is divisible by $p$. Also $c^p=c$ for every $c\in\F_p$ (Fermat's little theorem). So
$\bigl(\sum_nu_nx^n\bigr)^p=\sum_nu_n^px^{pn}=\sum_nu_nx^{pn}=U(x^p)$. The first equality holds for polynomials by
repeated use of $(a+b)^p=a^p+b^p$; it holds for power series because each coefficient of $U(x)^p$ involves only finitely
many coefficients of $U$.
\end{proof}

For $p=2$ this says that squaring a series modulo $2$ doubles every exponent: for instance,
$(x+x^3)^2\equiv x^2+x^6\pmod2$. In particular, the square of any series has, modulo $2$, no terms of odd degree.

\section{Results}\label{sec:results}

\begin{theorem}[Exactly one integer solution]\label{thm:exist}
Equation \eqref{eq:oeis} has exactly one normalized solution $A$ with integer coefficients. Its compositional inverse
$R$ has integer coefficients and satisfies
\[
R(3x^2)=2x^2+R(x)^2+2R(x)^3 .
\]
\end{theorem}
(Applying $A$ to both sides of the last equation gives formula (E2) of the entry.)

\begin{theorem}[The square modulo $9$]\label{thm:nine}
$A(x)^2\equiv x^2+2x^3\pmod 9$. In particular, $\coeff{n}{A(x)^2}\equiv0\pmod 9$ for every $n>3$.
\end{theorem}

\begin{proposition}[A classical identity]\label{prop:frob}
For every power series $W$ with integer coefficients and all integers $k,n\ge0$,
\[
\coeff{n}{\bigl(W(x)^{3k}-W(x^3)^k\bigr)}\equiv0\pmod 3 .
\]
\end{proposition}

\begin{corollary}[Hanna's conjectures]\label{cor:conj}
Statements (E3) and (E4) of the entry hold: $\coeff{n}{A(x)^2}\equiv0\pmod9$ for $n>3$, and
$\coeff{n}{\bigl(A(x)^{3k}-A(x^3)^k\bigr)}\equiv0\pmod3$ for all $n,k\ge1$ (indeed for all $n,k\ge0$).
\end{corollary}
Corollary~\ref{cor:conj} is immediate from Theorem~\ref{thm:nine} and from Proposition~\ref{prop:frob} with $W=A$.

\section{Proofs}\label{sec:proofs}

\para{The idea in brief.} We first study the inverse series $R$. It satisfies the equation $R(3x^2)=2x^2+R^2+2R^3$, in which the only
substitution is that of the monomial $3x^2$ into the unknown $R$. Lemma~\ref{lem:family} shows that this equation has exactly one integer
solution, by checking that the division by $2$ needed at each step is always exact. Its inverse is $A$
(Theorem~\ref{thm:exist}). Theorem~\ref{thm:nine} then follows by applying $R$ to \eqref{eq:oeis} and working
modulo $9$, where $(3A^2)^2=9A^4$ vanishes.

\subsection{An equation with an exact halving step}
Both this paper and its companion paper on OEIS A389475 rest on the following lemma, which concerns the
equation
\begin{equation}\label{eq:family}
R(3x^2)=2x^2+R(x)^2+2s\,R(x)^3 ,
\end{equation}
where $s$ is a fixed integer. Here $R(3x^2)$ is $R$ with $x$ replaced by $3x^2$, so its coefficient of $x^{2k}$ is
$3^kr_k$, and $R(x)^2$, $R(x)^3$ are ordinary powers.

\emph{Why an argument is needed.} Suppose we compute the coefficients $r_2,r_3,\dots$ of a normalized
solution $R=x+r_2x^2+r_3x^3+\cdots$ one at a time. The coefficient $r_n$ first appears in degree $n+1$, in the square
$R(x)^2=x^2+\cdots+2r_nx^{n+1}+\cdots$, with the factor $2$. (In $R(x)^3$ it appears only from degree $n+2$ on, and in
$R(3x^2)$ only in degree $2n$.) So each step requires a division by $2$, and it is not clear that the result is an
integer. The proof shows that the number to be divided is always even.

\begin{lemma}[Exact halving]\label{lem:family}
For every integer $s$ there is exactly one normalized series $R_s$ with integer coefficients that satisfies
\eqref{eq:family}.
\end{lemma}

\begin{proof}
For a normalized series $U$ with integer coefficients, define its \emph{error}
\[
E(U)=U(x)^2+2s\,U(x)^3-U(3x^2)+2x^2 ,
\]
so that $U$ solves \eqref{eq:family} exactly when $E(U)=0$.

\emph{Step 1: the error is always even.} Reduce modulo $2$. By Lemma~\ref{lem:frobenius} with $p=2$,
$U(x)^2\equiv U(x^2)$. Since $3\equiv1$, also $U(3x^2)\equiv U(x^2)$. The terms $2s\,U^3$ and $2x^2$ vanish modulo $2$.
So $E(U)\equiv0\pmod2$: every coefficient of $E(U)$ is even. Moreover $E(U)$ has no terms of degree $0$, $1$ or $2$:
in degree $2$ its coefficient is $1-3+2=0$.

\emph{Step 2: changing one coefficient.} Let $V$ be another normalized integer series with
$V-U=c\,x^r+O(x^{r+1})$, where $r\ge2$. Then
\begin{itemize}
\item $V^2-U^2=(V-U)(V+U)=2c\,x^{r+1}+O(x^{r+2})$, because $V+U=2x+O(x^2)$;
\item $V^3-U^3=(V-U)(V^2+VU+U^2)=O(x^{r+2})$, because $V^2+VU+U^2=3x^2+O(x^3)$;
\item $V(3x^2)-U(3x^2)=O(x^{2r})$, which is $O(x^{r+2})$ because $r\ge2$.
\end{itemize}
Hence $E(V)-E(U)=2c\,x^{r+1}+O(x^{r+2})$. The same computation shows that if $V-U=O(x^{r+1})$, then
$E(V)-E(U)=O(x^{r+2})$.

\emph{Step 3: construction.} Start with $U=x$, and suppose that $E(U)=O(x^{r+1})$ for some $r\ge2$ (this holds for
$r=2$ by Step~1). Let $e$ be the coefficient of $x^{r+1}$ in $E(U)$. By Step~1, $e$ is even, so $e/2$ is an integer.
Replace $U$ by $U-(e/2)\,x^r$. By Step~2 the new error is $E(U)-e\,x^{r+1}+O(x^{r+2})=O(x^{r+2})$, and the coefficients
of $U$ below degree $r$ have not changed. Repeating this for $r=2,3,4,\dots$ fixes the coefficients of $U$ one at a
time. Let $U_r$ denote the series obtained after the step for degree $r$, so that $E(U_r)=O(x^{r+2})$. The
resulting series $R_s$ has integer coefficients, and $R_s-U_r=O(x^{r+1})$ for every $r$. By Step~2,
$E(R_s)-E(U_r)=O(x^{r+2})$, so $E(R_s)=O(x^{r+2})$ for every $r$; that is, $E(R_s)=0$.

\emph{Step 4: uniqueness.} Suppose two different normalized solutions first differ in degree $r$, by $c\ne0$. Then
$r\ge2$, and by Step~2 their errors differ by $2c\ne0$ in degree $r+1$. But both errors are $0$, a contradiction.
\end{proof}

\subsection{Proofs of the results}

\begin{proof}[Proof of Theorem~\ref{thm:exist}]
\emph{Existence.} Let $R=R_1$ be the series of Lemma~\ref{lem:family} with $s=1$, and let $A$ be its compositional
inverse (Lemma~\ref{lem:inverse}); both have integer coefficients. Substitute $A(x)$ for $x$ in \eqref{eq:family}.
Since $R(A(x))=x$, this gives
\[
R\bigl(3A(x)^2\bigr)=2A(x)^2+x^2+2x^3 .
\]
Substituting both sides into $A$, and using $A(R(y))=y$, gives $3A(x)^2=A\bigl(x^2+2x^3+2A(x)^2\bigr)$, which is
\eqref{eq:oeis}.

\emph{Uniqueness.} Let $U$ be any normalized integer solution of \eqref{eq:oeis}, and let $V$ be its compositional
inverse. Substituting $V(x)$ for $x$ in \eqref{eq:oeis}, and using $U(V(x))=x$, gives
$U\bigl(V(x)^2+2V(x)^3+2x^2\bigr)=3x^2$. Substituting both sides into $V$ gives $V(x)^2+2V(x)^3+2x^2=V(3x^2)$. So $V$
solves \eqref{eq:family} with $s=1$, hence $V=R_1$ by Lemma~\ref{lem:family}, and $U$, the inverse of $V$, equals $A$.
\end{proof}

\begin{proof}[Proof of Theorem~\ref{thm:nine}]
Work with coefficients modulo $9$, and write $S=A(x)^2$ and $Q=x^2+2x^3$, so that \eqref{eq:oeis} reads
$A(Q+2S)=3S$. Let $R$ be the inverse of $A$ (reduced modulo $9$). Substituting both sides into $R$ gives
\[
Q+2S=R(3S).
\]
Now $R(t)=t+t^2H(t)$ for some series $H$, so $R(3S)=3S+(3S)^2H(3S)=3S+9S^2H(3S)=3S$, because $9=0$ modulo $9$. Hence
$Q+2S=3S$, that is, $S=Q$. No division by $3$ was used.
\end{proof}

\begin{proof}[Proof of Proposition~\ref{prop:frob}]
Reduce modulo $3$. By Lemma~\ref{lem:frobenius} with $p=3$, $W(x)^3\equiv W(x^3)$. Raising both sides to the $k$-th
power gives $W(x)^{3k}\equiv W(x^3)^k\pmod 3$.
\end{proof}

\section{Checking the results}\label{sec:checks}

A mathematical proof can contain mistakes, so the results were checked in two independent ways.

\para{Numerical check.}
We solved \eqref{eq:oeis} coefficient by coefficient with exact rational arithmetic, up to degree $89$. All
coefficients came out as integers, and they reproduce all 24 terms printed in the OEIS entry. Theorem~\ref{thm:nine}
holds for every coefficient of degree below $90$, and statement (E4) for $k\le6$ and every coefficient of degree below
$90$.
The computation was done with the PARI/GP computer algebra system, using code that was written separately from the
formal proofs (file \texttt{verification/independent\_checks.gp} in the accompanying repository~\cite{repo}).
A computation can only test finitely many cases, so it is a check on the definitions and on the indexing, not a proof.

\para{Formal proof in Lean.}
Lean~\cite{lean} is a \emph{proof assistant}: a program in which mathematical statements and their proofs are written
in a precise formal language, and every step is checked mechanically down to the basic rules of logic. The theorems of Section~\ref{sec:results} are proved in Lean~4 (version 4.19.0) on top of
Mathlib~\cite{mathlib}, Lean's library of formalized mathematics. The statements are collected in the file \texttt{lean/Entries/A389474.lean} of the repository,
in the namespace \texttt{OEIS.A389474}. Lemma~\ref{lem:family} is \texttt{integral\_family\_existence}. The constructions are in \texttt{lean/Families/CubicReversion.lean}, and the congruences use \texttt{lean/Families/CubicFrobenius.lean} and \texttt{lean/Families/CubicConsequences.lean}.
The table lists where each result of this paper is proved.
\begin{center}\small\renewcommand{\arraystretch}{1.25}
\begin{tabular}{>{\raggedright\arraybackslash}p{0.52\linewidth}>{\raggedright\arraybackslash}p{0.4\linewidth}}
\toprule
Statement in this paper & Lean theorem(s)\\
\midrule
Lemma~\ref{lem:family}, for every integer $s$ & \texttt{integral\_\allowbreak{}family\_\allowbreak{}existence} \\
Theorem~\ref{thm:exist}: $A$ is normalized, solves \eqref{eq:oeis}, and is the only normalized integer solution & \texttt{normalized}, \texttt{source\_\allowbreak{}equation}, \texttt{source\_\allowbreak{}unique}, \texttt{source\_\allowbreak{}existence\_\allowbreak{}uniqueness} \\
Theorem~\ref{thm:nine}, as a series and coefficient by coefficient & \texttt{square\_\allowbreak{}mod\_\allowbreak{}nine}, \texttt{square\_\allowbreak{}coefficients} \\
Proposition~\ref{prop:frob} for $W=A$ (all $n,k\ge0$) & \texttt{power\_\allowbreak{}coefficients} \\
\bottomrule
\end{tabular}
\end{center}
The Lean project was rebuilt from source, and the axioms on which each of these theorems depends were listed with
Lean's \texttt{\#print axioms} command. Every one depends only on the three standard axioms of Lean's logic
(\texttt{propext}, \texttt{Classical.choice}, \texttt{Quot.sound}); no step is left unproved. A formal proof certifies
exactly the statement that Lean checks. Whether that statement says what the OEIS entry means is a separate question,
so the Lean definitions were also compared with the definition in the OEIS entry
(Section~\ref{sec:sequence}).

\section{Discussion}\label{sec:discussion}

\para{Scope.} Corollary~\ref{cor:conj} settles both conjectures of the entry completely. Statement (E4) is not special to
this sequence: it holds for every power series with integer coefficients.

\para{The series itself modulo $9$.} Theorem~\ref{thm:nine} determines $A$ itself modulo $9$:
\[
A(x)\equiv x\,(1+2x)^{1/2}=x+x^2-\tfrac12x^3+\tfrac12x^4-\tfrac58x^5+\cdots\pmod 9 .
\]
Here $(1+2x)^{1/2}=\sum_n\binom{1/2}{n}(2x)^n$ is the binomial series. Its coefficients are fractions whose
denominators are powers of $2$, and such fractions make sense modulo $9$ because $2$ has an inverse there
($2\cdot5=10\equiv1$, so $\frac12\equiv5$). For example, $-\frac12\equiv4\equiv a(3)$ and $\frac12\equiv5\equiv a(4)=14$.
\emph{Proof.} Modulo $9$, the series $W=A(x)/x$ has constant term $1$ and, by Theorem~\ref{thm:nine}, $W^2=1+2x$. The
binomial series $Y=(1+2x)^{1/2}$ also has constant term $1$ and $Y^2=1+2x$. Then $(W-Y)(W+Y)=0$, and $W+Y$ has
constant term $2$, which is invertible modulo $9$; so $W+Y$ has a multiplicative inverse, and $W=Y$. We checked this
formula by computer for every coefficient of degree below $90$. It is not part of the Lean formalization.

\para{Higher powers of $3$.} The square of $A$ remains a polynomial modulo every power of $3$. For example,
\[
A(x)^2\equiv x^2+2x^3+9x^4+9x^5+9x^6\pmod{27}.
\]
More precisely: \emph{for every $k\ge2$, the series $A(x)^2$ is congruent modulo $3^k$ to a polynomial of degree at most
$3k-3$.} \emph{Proof.} Keep the notation $S=A(x)^2$, $Q=x^2+2x^3$ from the proof of Theorem~\ref{thm:nine}, but now
work over the integers, where $Q+2S=R(3S)$ holds exactly. Put $K(u)=3R(u)-2u$, a normalized series with integer
coefficients. Then $K(3S)=3R(3S)-6S=3Q$, so $3S=\bar K(3Q)$, where $\bar K$ denotes the compositional inverse of $K$
(Lemma~\ref{lem:inverse}). Since $K(u)=u+3\bigl(R(u)-u\bigr)\equiv u\pmod 3$, also $\bar K(v)\equiv v\pmod3$, so
$\bar K(v)=v+\sum_{j\ge2}3c_jv^j$ with integers $c_j$. Therefore
$3S=3Q+\sum_{j\ge2}3^{j+1}c_jQ^j$, and dividing by $3$,
\[
A(x)^2=S=Q+\sum_{j\ge2}3^{\,j}c_j\,Q^j .
\]
Modulo $3^k$ only the terms with $j<k$ survive, which gives a polynomial in $Q$ of degree at most $k-1$, that is, a
polynomial in $x$ of degree at most $3k-3$. For $k=2$ this is Theorem~\ref{thm:nine} again. For $k=3$ one finds
$c_2=1$, which gives the formula modulo $27$ above. We checked the statement by computer for $k\le7$ and every coefficient
of degree below $90$. It is not part of the Lean formalization.

\para{Open questions.} The method does not say much about the coefficients $a(n)$ of $A$ modulo higher powers of $3$,
or modulo other primes. The companion entry A389475 is a second member of the family of Lemma~\ref{lem:family} (the case
$s=-1$); it would be interesting to understand the other members.

\section*{About this paper}\label{sec:about}
\para{How the work was done.}
This paper is one of 29 papers on conjectures from the OEIS that came out of an AI research effort directed by
Dev~D.~Goyal~\cite{repo}. The roles were as follows.
\begin{itemize}
\item \emph{Dev~D.~Goyal} started the project, chose which problems to pursue, asked for proofs, formal verification and
revisions, and decided what to publish. Dev~D.~Goyal did not supply proof strategies or intermediate arguments.
\item \emph{ChatGPT} (OpenAI), using the GPT-6 model with Pro reasoning, worked autonomously from those prompts. It
searched the OEIS for suitable conjectures, found the proofs, wrote the Lean~4 formalization, and wrote a first draft
of this paper.
\item \emph{Claude Opus~5.5}, used with Claude Code, checked the work independently. It rebuilt the Lean project from source and confirmed that the theorems listed in Section~\ref{sec:checks} rely only on
Lean's standard axioms; compared the Lean definitions with the OEIS definition; recomputed the sequence with separately
written PARI/GP code, compared it with the terms published in the OEIS, and tested every claim of the paper on many
terms; and read every step of the proofs. It then wrote the present version of the paper, expanding the
explanations and, in places, writing out an argument in more detail or by a more elementary route. The theorems of
Section~\ref{sec:results} and the proof strategies are ChatGPT's. The Discussion section, including the short arguments
in it, was written by Claude Opus~5.5; its observations were checked by computer and are not part of the Lean
formalization.
\end{itemize}
Dev~D.~Goyal has verified the results and is responsible for this paper.

\para{Review status.} This paper has not been peer reviewed, and it has not been submitted to a journal or to the OEIS.

\para{Priority.} An OEIS entry that labels a statement as a conjecture shows only that its contributor did not know
a proof; it does not show that no proof existed. During the research only limited searches for earlier work were
made. No claim of priority is made.

\para{Corrections.} Please report any error in the issue tracker of the repository~\cite{repo}, where the Lean sources and the numerical checks are also
available.

\section*{Acknowledgements}
The sequence and the conjectures are due to Paul~D.~Hanna~\cite{oeis}, who also contributed the companion entry
A389475~\cite{oeisother}. Statement (E4) is an instance of the classical Frobenius identity and is not claimed as new.

\begin{thebibliography}{99}\small
\raggedright
\bibitem{oeis} P.~D.~Hanna, entry A389474 (G.f.\ $A(x)$ satisfies $A(x^2+2x^3+2A(x)^2)=3A(x)^2$), \emph{The On-Line Encyclopedia of Integer Sequences}, October 2025. \url{https://oeis.org/A389474}. Text as of October 2026.
\bibitem{oeisother} P.~D.~Hanna, entry A389475 (G.f.\ $A(x)$ satisfies $A(3x^2)=2x^2+A(x)^2-2A(x)^3$), \emph{The On-Line Encyclopedia of Integer Sequences}, October 2025. \url{https://oeis.org/A389475}. Text as of October 2026.
\bibitem{oeisdb} N.~J.~A.~Sloane and The OEIS Foundation Inc., \emph{The On-Line Encyclopedia of Integer Sequences}. \url{https://oeis.org}.
\bibitem{stanley} R.~P.~Stanley, \emph{Enumerative Combinatorics}, Volume~2, Cambridge University Press, 1999. (Chapter~5 treats formal power series, composition and compositional inverses.)
\bibitem{niven} I.~Niven, Formal power series, \emph{American Mathematical Monthly} 76 (1969), 871--889.
\bibitem{lean} L.~de~Moura and S.~Ullrich, The Lean~4 theorem prover and programming language, in \emph{Automated Deduction -- CADE 28}, Lecture Notes in Computer Science 12699, Springer, 2021, 625--635.
\bibitem{mathlib} The mathlib Community, The Lean mathematical library, in \emph{Proceedings of the 9th ACM SIGPLAN International Conference on Certified Programs and Proofs (CPP 2020)}, 367--381. Version used: commit \texttt{c44e0c8ee63c}.
\bibitem{repo} D.~D.~Goyal, \emph{Proofs and corrections of OEIS conjectures}, repository, 2026. \url{https://github.com/dvgyl/oeis-a396843-proof}.
\end{thebibliography}
\end{document}
